% Teaching manuscript. Edit this TeX source; Markdown is compatibility navigation only.


\StudyChapter{CS5489-Tutorial01}{Tutorial 1 · From Python loops to orthonormal vectors}
\StudyTarget{CS5489-Tutorial01__cell-1}
\StudyTarget{CS5489-Tutorial01__notebook-1f98a61d}
\StudyTarget{CS5489-Tutorial01__tutorial-1-from-python-loops-to-orthonormal-vectors}



\StudyHeading{1}{从变量变化，走到能验证的程序}
\StudyTarget{CS5489-Tutorial01___1}


\StudyRef{CS5489-Lecture01}{root}{Lecture 1} · \StudyRef{Foundations-MathForML}{projection}{数学补课} · \StudyRepoRef{CS5489/course-notes/README.md}{}{课程目录}





原任务顺序：认识Notebook状态 \ensuremath{\to} 数非平凡因子 \ensuremath{\to} 折线图和直方图 \ensuremath{\to} 读取数据 \ensuremath{\to} Gram–Schmidt。先跟随下方示范理解每步的输入与输出，再在配套Notebook中从头运行，最后合上答案自己写一次。

依据当前原Tutorial1的47个单元编写；单元号从1起算，不是In[n]。本讲义解释配套Notebook的代码与已保存输出，例题答案是整理者实现。原课的Colab挂盘改为本地文件路径，练习输出写入独立目录。


\begingroup
\small\setlength{\tabcolsep}{3pt}
\renewcommand{\arraystretch}{1.22}
\begin{longtable}{@{}>{\raggedright\arraybackslash}p{\dimexpr0.374772\linewidth-1.50000\tabcolsep\relax}>{\raggedright\arraybackslash}p{\dimexpr0.163207\linewidth-1.50000\tabcolsep\relax}>{\raggedright\arraybackslash}p{\dimexpr0.203980\linewidth-1.50000\tabcolsep\relax}>{\raggedright\arraybackslash}p{\dimexpr0.258041\linewidth-1.50000\tabcolsep\relax}@{}}
\toprule
原任务 & 优先级 & 难度与卡点 & 掌握要求与依据 \\
\midrule\endfirsthead
原任务 & 优先级 & 难度与卡点 & 掌握要求与依据 \\
\midrule\endhead
\bottomrule\endfoot
Tutorial1，第1–13个单元 Notebook与变量 & 核心必会 & 低，执行顺序 & 能从干净内核复现；当前原任务 \\
Tutorial1，第14–25个单元 因子与图 & 核心必会 & 中，循环范围和两种横轴 & 能独立实现并解释输出 \\
Tutorial1，第26–34个单元 数据文件 & 常规掌握 & 中，路径与数组轴 & 能读取、核对shape、画点 \\
Tutorial1，第35–45个单元 Gram–Schmidt & 核心必会 & 高，投影与张成空间 & 实现并检查正交、长度与空间保持 \\
Tutorial1，第46–47个单元 查帮助 & 常规掌握 & 低，工具用法 & 会查接口而非猜参数 \\
\end{longtable}
\endgroup






\StudyHeading{1}{0. Setup｜本地路径与输出位置}
\StudyTarget{CS5489-Tutorial01__cell-2}
\StudyTarget{CS5489-Tutorial01__notebook-46e80bb2}
\StudyTarget{CS5489-Tutorial01__0-setup}


只在本节修改路径。例题新文件和图片放在本讲义输出目录，原件只读。
English: Configure paths here; generated files go to the learning-note output directory.




\StudyTarget{CS5489-Tutorial01__cell-3}
\StudyTarget{CS5489-Tutorial01__notebook-363897ef}
\StudyNotebookCell{CS5489/course-notes/Tutorial01.ipynb}{363897ef}





\StudyHeading{1}{1. Notebook state｜屏幕顺序不等于执行顺序}
\StudyTarget{CS5489-Tutorial01__cell-4}
\StudyTarget{CS5489-Tutorial01__notebook-8328b3fc}
\StudyTarget{CS5489-Tutorial01__1-notebook-state}


\textbf{原任务 / Task：} 运行单元、重复更新变量、重启内核观察状态。 / Run cells, repeat an update, and understand kernel restart.


\StudyInline{x=\allowbreak{}x+1}先读取旧x，再算右侧，最后把新值赋给x。以下把三次运行展开，保存清楚的状态轨迹。示例捕获NameError，便于继续观察后续单元。


来源：Tutorial1，第3–13个单元。




\StudyTarget{CS5489-Tutorial01__cell-5}
\StudyTarget{CS5489-Tutorial01__notebook-d5030717}
\StudyNotebookCell{CS5489/course-notes/Tutorial01.ipynb}{d5030717}




\StudyTarget{CS5489-Tutorial01__cell-6}
\StudyTarget{CS5489-Tutorial01__notebook-f523b258}
\textbf{独立题 / Try:} 初始x=10，把\StudyInline{x=\allowbreak{}x*2}执行两次，得到多少？重启内核后直接运行这一行会怎样？ / Starting at x=10, execute x=x*2 twice. What happens if you restart the kernel and run only that line?


\textbf{答 / Answer:} 40；重启清除了变量，需重新定义x。 / 40; after restart x must be initialized again.





\StudyHeading{1}{2. Count factors｜循环范围与因子计数}
\StudyTarget{CS5489-Tutorial01__cell-7}
\StudyTarget{CS5489-Tutorial01__notebook-1a9645aa}
\StudyTarget{CS5489-Tutorial01__2-count-factors}


\textbf{原任务 / Task：} 对2到100的每个整数，统计除1和自身之外的正因子数。 / Count positive divisors excluding 1 and n for every integer n from 2 through 100.


\begin{StudyBlock}{跟做}{先用12检查循环与返回位置}
对n=12，逐个试d=2,…,11；\StudyInline{n \% d =\allowbreak{}=\allowbreak{} 0}表示除后余数为0。2、3、4、6通过，得到4。\StudyInline{return}应在整个循环结束后执行，不能发现第一个因子就离开。
\end{StudyBlock}


来源：Tutorial1，第14–18个单元。




\StudyTarget{CS5489-Tutorial01__cell-8}
\StudyTarget{CS5489-Tutorial01__notebook-88fd2593}
\StudyNotebookCell{CS5489/course-notes/Tutorial01.ipynb}{88fd2593}




\StudyTarget{CS5489-Tutorial01__cell-9}
\StudyTarget{CS5489-Tutorial01__notebook-43df6a5d}
\textbf{变式 / Transfer:} 不运行程序，求9和13的结果。 / Without running the program, find the counts for 9 and 13.
\textbf{答 / Answer:} 9除1和自身外只有正因子3，得到1；13为素数，得到0。 / The only positive divisor of 9 other than 1 and itself is 3, giving 1; 13 is prime, giving 0.


第一张图问“每个数有几个因子”，第二张图问“有多少数落在某个因子数里”。两个横轴不是同一个东西。




\StudyTarget{CS5489-Tutorial01__cell-10}
\StudyTarget{CS5489-Tutorial01__notebook-f3a5d03f}
\StudyNotebookCell{CS5489/course-notes/Tutorial01.ipynb}{f3a5d03f}




\StudyTarget{CS5489-Tutorial01__cell-11}
\StudyTarget{CS5489-Tutorial01__notebook-4baebf46}
\StudyNotebookCell{CS5489/course-notes/Tutorial01.ipynb}{4baebf46}




\StudyTarget{CS5489-Tutorial01__cell-12}
\StudyTarget{CS5489-Tutorial01__notebook-b611bf34}
看图检查：0这一箱对应素数，共25个；所有箱合计99。原固定坐标范围对应本题，换输入范围后应重新检查是否截掉数据。
English: The zero bin contains the 25 primes in this range. Histogram counts total 99; fixed axis limits need reconsideration when the range changes.





\StudyHeading{1}{3. Load and plot the supplied array｜先核对两根轴}
\StudyTarget{CS5489-Tutorial01__cell-13}
\StudyTarget{CS5489-Tutorial01__notebook-b1ed5995}
\StudyTarget{CS5489-Tutorial01__3-load-and-plot-the-supplied-array}


\textbf{原任务 / Task：} 从原data1.pickle读入mydata，检查shape，画第1列对第2列。 / Load mydata, inspect its shape, and plot column 1 against column 2.


下面读取课程提供的NumPy数组文件。Pickle加载可能执行对象重建操作，因此这里只使用已知来源的课程文件。


来源：Tutorial1，第26–34个单元。




\StudyTarget{CS5489-Tutorial01__cell-14}
\StudyTarget{CS5489-Tutorial01__notebook-8c598fb8}
\StudyNotebookCell{CS5489/course-notes/Tutorial01.ipynb}{8c598fb8}




\StudyTarget{CS5489-Tutorial01__cell-15}
\StudyTarget{CS5489-Tutorial01__notebook-30f72bf7}
\StudyInline{mydata[:\allowbreak{},\allowbreak{}0]}选所有行的第0列，不是第一行。数据形状为120\ensuremath{\times}2；原文件未给出类别标签或坐标的物理单位。
\textbf{自查 / Check:} \StudyInline{mydata[:\allowbreak{}3,\allowbreak{}:\allowbreak{}]}取什么？ / What does mydata[:3,:] select?
\textbf{答 / Answer:} 前三条样本的全部特征，形状(3,2)。 / All features of the first three samples, shape (3,2).




\StudyTarget{CS5489-Tutorial01__cell-16}
\StudyTarget{CS5489-Tutorial01__notebook-2352c6a8}




\StudyHeading{1}{4. Gram–Schmidt｜拿走已有方向，剩下才是新信息}
\StudyTarget{CS5489-Tutorial01__gram-schmidt}
\StudyTarget{CS5489-Tutorial01__4-gramschmidt}


\textbf{原任务 / Task：} 输入矩阵V每列是一个向量，返回同一张成空间的正交归一列E。 / Return orthonormal columns spanning the same column space as V.


先补\StudyRef{Foundations-MathForML}{projection}{长度、投影与张成空间}。例如v=(2,1)、u=(1,0)，投影到u方向是(2,0)，相减余下(0,1)。不是把所有向量分别除以长度就能让它们互相垂直。


对非零u，\(P_u(v)=(v^Tu)/(u^Tu)\,u\)。原步骤是从第k列减去已有各方向的投影，再归一化。下面存储单位方向q，因此投影简化为\((q^Tv)q\)。

\begin{StudyBlock}{完整示范 / Worked example}{两列向量的正交化}
把三步用一对具体列走完（教学补充）：\(v_1=(1,1)^T,v_2=(1,0)^T\)。第一列长度是\(\sqrt2\)，所以\(e_1=(1,1)^T/\sqrt2\)；第二列在它上面的投影是\((e_1^Tv_2)e_1=(1/2,1/2)^T\)。减去投影，余量\((1/2,-1/2)^T\)的长度是\(1/\sqrt2\)，归一化得到\(e_2=(1,-1)^T/\sqrt2\)。两列各自长度1，内积0，且能重建原来两列。


\textbf{English worked example:} For columns (1,1)\textsuperscript{T} and (1,0)\textsuperscript{T}, normalize the first, subtract the second column's projection (1/2,1/2)\textsuperscript{T}, then normalize its residual. The resulting columns are (1,1)\textsuperscript{T}/\(\sqrt{2}\) and (1,\ensuremath{-}1)\textsuperscript{T}/\(\sqrt{2}\).
\end{StudyBlock}

这对应原课的classical Gram–Schmidt。数值接近线性相关时会更难稳定计算，后面的检查会帮助发现问题；其他QR算法可在二读时学习。


来源：Tutorial1，第35–45个单元。




\StudyTarget{CS5489-Tutorial01__cell-17}
\StudyTarget{CS5489-Tutorial01__notebook-a6816ae8}
\StudyNotebookCell{CS5489/course-notes/Tutorial01.ipynb}{a6816ae8}




\StudyTarget{CS5489-Tutorial01__cell-18}
\StudyTarget{CS5489-Tutorial01__notebook-94a085a9}
\StudyNotebookCell{CS5489/course-notes/Tutorial01.ipynb}{94a085a9}




\StudyTarget{CS5489-Tutorial01__cell-19}
\StudyTarget{CS5489-Tutorial01__notebook-2d827844}
\begin{StudyBlock}{输出解读 / Interpretation}{长度、正交与张成空间}
对角线1表示单位长度，非对角线接近0表示正交；另一项重建检查才说明没有换成无关的子空间。 / Diagonal ones check unit lengths; near-zero off-diagonal entries check orthogonality. Reconstruction checks that the original span was retained.
\end{StudyBlock}


\textbf{独立变式 / Transfer:} V的两列是(1,0)和(1,1)。手算E；如果第二列改成(2,0)又怎样？ / Hand-compute E for columns (1,0) and (1,1). What changes if the second column becomes (2,0)?
\textbf{答 / Answer:} 得到(1,0)、(0,1)；改后第二列残差为零，只提供一个方向，不能产生两个正交单位列。 / The basis is (1,0),(0,1). The replacement is dependent and yields a zero residual, so two unit directions cannot be produced.




\StudyTarget{CS5489-Tutorial01__cell-20}
\StudyTarget{CS5489-Tutorial01__notebook-129ae6f4}
\StudyNotebookCell{CS5489/course-notes/Tutorial01.ipynb}{129ae6f4}





\StudyHeading{1}{5. Help and return｜最后练一次英文解释}
\StudyTarget{CS5489-Tutorial01__cell-21}
\StudyTarget{CS5489-Tutorial01__notebook-2f911c2d}
\StudyTarget{CS5489-Tutorial01__5-help-and-return}


查询函数用法时，可用\StudyInline{help(\allowbreak{}myGS)\allowbreak{}}或Jupyter中的\StudyInline{myGS?}查看接口。来源：Tutorial1，第46–47个单元。调试先读异常最后一行和对应代码位置；重启并从头运行，能排除“旧变量碰巧还在”的情况。


\textbf{Q / 题：} Why is E\textsuperscript{T}E=I not the only check? / 为什么只检查E\textsuperscript{T}E=I还不够？
\textbf{A / 答：} An unrelated orthonormal basis can also pass it; check that E reconstructs every original column. / 无关的正交基也能通过，还要检查原列能由E重建。


\StudyRef{CS5489-Lecture01}{linear-algebra}{回Lecture1} · \StudyRef{CS5489-Lecture02}{root}{下一讲Lecture2} · \href{https://crazyshout.github.io/micro-course/?lesson=ml04}{微课ML04}。按课程目录提供的已有卡号巩固，算法仍需独立重写。



\StudyHeading{1}{原文件覆盖索引}
\StudyTarget{CS5489-Tutorial01___1-2}


\begingroup
\small\setlength{\tabcolsep}{3pt}
\renewcommand{\arraystretch}{1.22}
\begin{longtable}{@{}>{\raggedright\arraybackslash}p{\dimexpr0.509633\linewidth-1.00000\tabcolsep\relax}>{\raggedright\arraybackslash}p{\dimexpr0.490367\linewidth-1.00000\tabcolsep\relax}@{}}
\toprule
原单元 & 本文件对应 \\
\midrule\endfirsthead
原单元 & 本文件对应 \\
\midrule\endhead
\bottomrule\endfoot
Tutorial1，第1–6个单元 & 标题、Setup；封面与运行设置；Colab改为本地路径 \\
Tutorial1，第7–13个单元 & §1状态轨迹与预期NameError \\
Tutorial1，第14–18个单元 & §2因子程序 \\
Tutorial1，第19–25个单元 & §2两幅原任务图 \\
Tutorial1，第26–34个单元 & §3数据读取、shape、散点 \\
Tutorial1，第35–45个单元 & §4投影、实现、原矩阵及三种检查 \\
Tutorial1，第46–47个单元 & §5帮助与收尾；Tutorial1，第47个单元为空单元 \\
\end{longtable}
\endgroup
